\documentclass[11pt]{article}
\usepackage{iftex}
\ifPDFTeX
  \usepackage[T1]{fontenc}
  \usepackage[utf8]{inputenc}
  \usepackage{tgpagella}
\else
  \usepackage{fontspec}
  \setmainfont{texgyrepagella}[Extension=.otf,UprightFont=*-regular,BoldFont=*-bold,ItalicFont=*-italic,BoldItalicFont=*-bolditalic]
\fi
\usepackage[margin=1.15in]{geometry}
\usepackage{setspace}
\usepackage{url}
\setstretch{1.08}
\setlength{\parindent}{1.4em}
\setlength{\parskip}{0.2em}

\title{Metered Memory:\\[0.3em]
\large Storage, Compute, and the Argument for Charging Nothing for Either}
\author{Flyxion\\Independent Researcher}
\date{October 2026}

\begin{document}
\maketitle

\begin{abstract}
\noindent Storage and compute are sold to the public by the month, by the gigabyte, by the second and by the token, at prices set by the handful of firms that dominate the market and enforced by the threat of deletion. This essay argues that both should instead be provided free at the point of use, up to a baseline sufficient for a person's lifelong record and ordinary computation, funded as infrastructure in the way roads, libraries and water mains are funded. The argument rests on five mechanisms. The price of stored data is not determined by the price of the disks that hold it. The charge for taking data out places the cost of leaving on the customer at the moment of leaving, and Europe has scheduled the end of such switching charges for January 2027. Free tiers gather the public's records and later make further deposits count against a finite allowance. Quotas count what the customer can see and not what the provider must store, so copies and revisions are billed by their visible volume. And hardware prices rise in ways the smallest buyers cannot negotiate, plausibly because the largest orders set the price. Each is set against published sources, and the limits of those sources are stated. Underneath the economics sits a normative argument that does not depend on them: a service that is competitively priced and efficiently run is still an unsuitable foundation for a lifelong record when continued access depends on individual purchasing power. The essay then argues that free capacity is not enough unless the means of rebuilding it are published and taught, says what free should mean, answers the strongest objections, and concedes where the case is incomplete.
\end{abstract}

\section{The Meter}

A meter turns a relationship into a quantity. Before the meter, a person who kept a shoebox of letters kept them; after it, the same person rents shelf space for the letters from a landlord who may raise the rent, change the terms, or close the account. The rent is small on any single day, which is how it works. The sums are set to be too small to fight and too regular to escape, and they accumulate for as long as the person wishes to remember anything.

The thesis of this essay is that the meter should come off. The premise of the essay is that storage, which is the capacity to keep a record, and compute, which is the capacity to transform one, have become the medium in which a person writes, argues, remembers, learns and is remembered. A society that charges its members by the gigabyte for the memory of their own lives, and by the token for the use of their own language, has made the basic conditions of literacy a subscription. The case against doing so is not sentimental. It rests on prices that can be read off public pages, on a European statute, on a corporate announcement, and on a vendor's own price table. The sections below set these out, then say what is proposed in their place.

A word on what is not claimed. The claim is not that the physical world is free. Disks cost money, current costs power, and buildings cost land. The claim is that a cost borne in common and recovered through taxation or cooperative levy is a different thing from a price charged to each person at each use by a seller who also decides the terms. Most roads cost money and are free at the point of use. Libraries cost money and lend without a till. The question is which kind of good a person's own record and a person's own computation should be, and the answer proposed here is that they belong with roads and libraries and not with taxi fares.

Two arguments run through the essay, and they bear different weight. The normative argument is that a person's lifelong record and everyday computation are the kind of capability a society should provide in common, whatever the economics. The political argument is that the present arrangement is not forced by cost, so that nothing in the economics stands in the way of the normative one. The empirical sections serve the second, and the normative argument, stated in full in the section on objections, would stand if they were weakened.

\section{Rent on a Disk}\label{sec:rent}

The first mechanism is the gap between what storage costs to buy and what it costs to rent. A published account of Amazon's S3 Standard storage gives the price as \$0.023 per gigabyte per month for the first fifty terabytes in the US East region \cite{filebase}. That is \$23 per terabyte per month. A price survey of new internal hard drives sold on Amazon in October 2026 puts the cheapest per-terabyte price across common capacities at roughly \$24 to \$41 \cite{hddhunt}. A month of renting a terabyte from the largest storage provider therefore costs about what it costs to own a terabyte of disk outright.

This is not a cost accounting, but it does bound the hardware component. The service replicates data across several facilities, so the customer is buying more than one disk, and the provider pays for power, buildings, staff and the engineering that makes the replication work. Standard storage is also the highest-priced tier of a range, and colder tiers are cheaper. Concede all of it. Triple the disk cost to allow for replication, which gives an illustrative hardware figure between \$72 and \$123 per terabyte. A customer who keeps that terabyte for one year pays \$276, between two and nearly four times that figure, and then pays again the next year and the year after, for as long as the file exists. Disks wear out and must be replaced over the life of a service, and tripling is only an illustration of replication, not an accounting of the other costs just listed, so none of this shows what the service costs to run or what margin it earns. What it shows is that the physical cost of the capacity does not determine the billing structure. The rent is perpetual, and it falls on the one thing that cannot be allowed to lapse, which is the record itself.

The point is made sharper by retrieval. The same survey of providers lists internet egress at \$0.090 per gigabyte for the first ten terabytes on AWS, \$0.087 on Azure and \$0.120 on Google Cloud at its premium tier, and \$0.00 on Cloudflare R2, which advertises unlimited free egress \cite{ecorpit}. At the AWS rate, reading one terabyte back out costs \$90, about four months of the \$23 monthly charge for storing it. Keeping a thing costs one amount and touching it costs another, larger one, and the second figure is the one that matters to a person who wants to move. The existence of a competitor that charges nothing for egress does not show that delivering bytes is free. It shows that billing for it separately is a commercial choice and not a law of physics, which is all the institutional argument needs.

\section{The Toll at the Exit}

The second mechanism is the exit fee. A fee for taking data out is not a fee for storing it, and its main effect is to make leaving expensive. The largest providers have acknowledged as much by retreat. According to the same survey, AWS, Azure and Google Cloud each waived egress fees in 2024 for customers fully leaving their platforms, and the waiver applies only to complete departures, not to routine transfers or the partial movement of a workload \cite{ecorpit}. A customer is thus permitted to leave if the customer leaves entirely and says so, and is charged if the customer merely wishes to keep a foot in another door.

Regulators have gone further. The European Data Act schedules the end of switching charges, including the egress charges incurred in switching, for 12 January 2027, with reduced charges permitted until then \cite{dataact}. A trade analysis reports that the ban covers infrastructure, platform and software services and binds non-European providers selling into the Union as well as European ones \cite{cloudmagazin}. The Regulation does not abolish every outbound transfer charge: its recitals allow providers to keep charging for egress used for in-parallel use of two services, though only up to the costs incurred. Nor does it ignore technical lock-in, since the recitals call for transfer in a commonly used, machine-readable format by means of open interfaces \cite{dataact}. The trade analysis nonetheless cautions that proprietary interfaces, vendor-specific formats and identity binding can still produce lock-in without any fee. The essay makes no claim about other jurisdictions, whose law is outside the sources used here.

The inference is narrower than it may appear. A charge that recovers a genuine cost is still a charge, and a genuine cost can still be a barrier to leaving. That the legislature assigned the cost of switching to the provider shows that cost recovery does not settle who should bear the cost when charging the customer impedes departure. The same reasoning supports the proposal below. A second lesson comes from the trade analysis: abolishing the toll does not by itself end lock-in, since format, interface and identity can remain the provider's to define. A baseline of free storage that is worth having must therefore come with open formats and portable identity as conditions of provision, and not as favors.

\section{The Free Tier}

The third mechanism is the free tier, which is how the public's records are gathered before they are billed. The pattern is visible from the beginning of the consumer cloud. Gmail was announced in 2004 with one gigabyte of storage, several hundred times what Hotmail and Yahoo offered at two or four megabytes \cite{sciencealert}. The offer worked as an argument as much as a product: keep everything, delete nothing, search instead of sorting. In May 2013 Google pooled Gmail, Drive and Google+ Photos under a single free allowance of 15 gigabytes, which raised the free total for many users, and extended paid storage to Gmail at the same time \cite{gmailblog}. That allowance was still 15 gigabytes when the next change arrived. On 11 November 2020 Google announced that the unlimited free storage of compressed photographs in Google Photos would end on 1 June 2021, after which new uploads would count against a free allowance of 15 gigabytes \cite{googlephotos}. Photographs already backed up under the old terms were exempted, and the announcement said the new allowance would hold, in its words, ``roughly three more years worth of memories'' for more than eighty percent of users. Beyond the allowance, storage was sold through a subscription plan priced in a Canadian report of the announcement at \$2.79 per month for one hundred gigabytes \cite{iphonecanada}.

Two features of this episode deserve attention, and neither is a charge of bad faith. The first is the exemption. The company honored the old uploads, which is to its credit, and the honoring shows that the arrangement was entirely a matter of discretion. A promise that holds for years and is then amended by announcement is a policy and not a right. The second is the sequence. An offer of unlimited free storage collects an enormous number of photographs. Once the collection exists and the habit of depositing has been formed, new deposits count against a finite allowance, with payment required beyond it. The people whose family records sit in the service are not the ones who chose the change, and they are invited to pay to go on adding to the archive past the allowance.

A record held at a provider's discretion is held conditionally, and conditions can change. The remedy is to move the record out of the category of things a firm gives away to attract customers and into the category of things a society provides because a life without them is diminished. A baseline allowance that cannot be withdrawn by announcement is the difference between a tenancy and a deed.

\section{Charging for Copies}

The fourth mechanism concerns what a quota counts. A quota counts what the user can see. A person who keeps a hundred revisions of a document, or who receives an attachment that ten thousand other people also received, is charged for a hundred documents and for one attachment per mailbox, whatever the provider does underneath. Whether the provider stores the bytes once or many times is the provider's business, and the user pays the same either way. Storage systems that identify content by a hash and keep a single copy of identical content, and that store a revision as the difference from its predecessor, are decades old and standard in version control. Where a provider uses such techniques internally, which the essay does not claim to know for any particular firm, the marginal cost of one more copy of identical content is close to nothing, and a quota that bills each copy as new is a price detached from the provider's cost by construction. It is a water company that bills every household for the whole reservoir because each household can see only its own tap.

The consequence is a tax on the habits that make a record valuable. Keeping drafts, keeping every revision and keeping what everyone else also keeps are the things a quota punishes. The rational response of the customer is to delete, to summarize or to export, which shrinks the record, and the record is what the customer came to the service to keep. A price that scaled with physical consumption would charge for what is new relative to what is already stored. How much that saves depends on the content. Identical files share storage entirely, and successive drafts of a text document differ by small changes that can be represented economically, but compressed photographs, video, encrypted data and heavily rewritten documents share little, and a long history of such material can cost far more than its final version. The essay promises no general ratio. The argument is that logical volume can differ from physical consumption, that the customer ordinarily cannot inspect the relationship, and that a price set on the logical figure alone leaves the whole difference with the seller. Billing by visible volume also teaches a habit, which is to treat history as a liability to be trimmed.

\section{The Price of Being Small}

The fifth mechanism is purchasing power. In September 2026 a server vendor, eRacks, published tracking figures for hard drive prices between March and August of that year \cite{eracks}. Its table gives cost per terabyte rising from \$16.25 to \$31.25 for one Seagate Barracuda line, from \$24.84 to \$41.22 for a WD Red Pro line, from \$21.78 to \$33.75 for a Seagate IronWolf Pro line, from \$29.79 to \$44.17 for a WD Gold line, from \$22.17 to \$30.64 for a Toshiba enterprise line and from \$24.33 to \$27.92 for a Seagate Exos line. The increases range from fifteen to ninety-two percent. The company attributes them to demand, of the same kind that raised memory and graphics-processor prices earlier in the year, and observes that consumer lines rose more than enterprise lines.

The source needs handling with care. It is a press release from a vendor describing its own price tracking, not an independent market study, and its headline range of forty to ninety percent does not match its own table, which runs from fifteen to ninety-two. The release names demand and does not name the buyers. That the buyers who set the price are the largest ones is therefore the essay's inference and not the source's finding, though it is not an eccentric inference: when a few purchasers can contract capacity in bulk and in advance, the ones who buy a drive at a time buy what is left, at the price the bulk buyers have established. The release also reports that the enterprise lines moved least and the consumer lines most, which complicates the inference. The pattern fits enterprise buyers having contracted ahead while spot buyers absorbed the increase, and it fits equally well consumer demand rising independently, and the source cannot distinguish them.

This bears on the thesis in two ways. It contradicts the easy slogan that storage is cheap, and the essay does not rest on that slogan. Whether a terabyte of disk costs sixteen dollars or forty-four, the price paid by a customer of a storage service is set by a separate decision. And it shows what an unprotected market does to everyone who is not a hyperscaler. The household that wants a disk, the library that wants a server and the researcher who wants a workstation do not negotiate. They take the price. A common baseline, bought at the scale of the purchasers who now set the terms, would put the public on the same side of the table as the firms that currently dominate the order book.

\section{Compute by the Token}

Compute is metered more finely than storage. A computation is billed by the second of processor time, by the hour of an accelerator, or by the token in a language-model interface, and each unit is small enough that no single charge invites a complaint. The fineness is the instrument. A price that is paid thousands of times a day is never examined, and a meter that runs at the speed of thought makes thinking itself an itemized expense. This section argues from the structure of public price lists, which anyone can read, and does not survey them.

The token shows the same detachment of price from work that the quota shows. It is a unit chosen by the seller for billing and not a measure of the computation done. Where providers discount repeated input, the discount concedes that some repetition is cheaper to serve, though a small change to a request can invalidate reusable computation, and a repeated prefix does not make a whole interaction cheap. Pricing by novelty would be no remedy, since it would require a definition of novelty and comparisons between users' requests that no one should want. The argument does not rest on any such meter. It rests on access, on predictable provision and on the value of iteration: most intellectual work proceeds by revising, re-asking and thinking again, and a meter that runs at each pass makes the person doing the work bear the price of every pass, at a rate the seller sets.

The case for treating compute like storage is, if anything, stronger, because compute is the more rapidly spreading medium of ordinary intellectual work. Searching a body of text, translating a letter, checking an argument, summarizing a report and drafting a reply are now done by machine, in the same way that arithmetic was once done on paper, and a person who cannot afford the machine is in the position of a person who cannot afford paper. Charging by the unit means that the cost falls hardest on the most active users, who are students, small researchers and anyone working on a long project, and that the rate is whatever the seller chooses. A customer cannot compare the price of a token against the cost of producing it, because that cost is not published.

None of this denies that large-scale computation is expensive. A training run at the frontier costs what it costs. The claim concerns the baseline: the computation an individual needs to read, write, search, calculate and keep a record is small by the standard of modern hardware, and the machinery that supplies it, whether in a public facility, a cooperative or the person's own device, can be provided at a flat level without metering. Above the baseline, supply can be priced at what it costs to provide it. The proposal is a floor and not a ceiling.

\section{Capacity Without Competence}

Free capacity delivered as an opaque service would repeat the enclosure under a different name. A person given an allotment by an operator whose storage works by mechanisms the person cannot inspect, in formats the person cannot read, behind an interface that explains nothing, holds a result and not the means of producing it. The distinction matters because the lock-in that survives the abolition of exit fees is of exactly this kind: proprietary interfaces, vendor-specific formats and identity binding remain in force after a ban on charges \cite{cloudmagazin}. What a society keeps, when it keeps anything, is the ability to rebuild the thing, and that ability lives in people and not in archives. Storage whose workings are published and taught leaves a population able to run its own, and storage whose workings are withheld leaves a population dependent on the operator's continued goodwill, however cheap the service.

Deduplication shows the difference. A provider that stores identical content once saves real money, and the saving can be kept as margin, passed on as a lower price, or explained to the public so that anyone can build the same system. The stored bytes are the same in all three cases. What differs is who can reproduce the efficiency afterward. Computation follows the same pattern. An assistant that supplies answers and an assistant that shows how its answers were reached may be equally accurate, and the population served by the second ends up more capable than the population served by the first. A baseline of free capacity should therefore be provided on terms that favor the second kind of system, with documented mechanisms, inspectable behavior and interfaces that teach.

\section{What Free Means}\label{sec:freemeans}

The word is used in a specific sense. It means free at the point of use for a defined baseline, provided as infrastructure and recovered through taxation or levy, in the way public libraries lend books without a till or public roads carry traffic without a toll. It does not mean that no resources are consumed, or that anyone may consume without limit. The analogy to roads and libraries concerns why a good is provided in common and not how economists classify it. Storage and compute are consumed in aggregate, as roads are congested, and the case for common provision rests on the same ground: exclusion cuts a person off from ordinary life. The provision has four conditions, and a scheme lacking any of them does not qualify.

The first condition is a baseline allotment, large enough to hold a person's lifetime of documents, photographs, correspondence and work, and to run the everyday computations of a literate life, set as a quantity and not as a promotion. Since a lifetime of video can exceed any fixed figure, the baseline is defined by principle and then fixed as a number: sufficient for documents, correspondence and still images across a lifetime, with bulk video above the baseline unless the number is raised, and reviewed periodically as hardware and habits change. An illustrative figure is on the order of a terabyte. The second is portability: open formats, published interfaces, no charge for export and an account identity the holder controls, so that the allowance can be moved to a different operator without penalty. The third is non-revocability within the baseline. A person's stored record cannot be deleted for non-payment, since there is no payment, and cannot be suspended at an operator's discretion; the sanction for abuse is a narrow one, applied through a procedure the holder can contest, and does not extend to the loss of the archive. The fourth condition is legibility: the mechanisms by which the capacity is provided are published and documented, so that the allowance can be understood and rebuilt as well as used.

Delivery can take several forms, and the argument does not require one. Operators can be public bodies, cooperatives, or private firms paid by a public authority to supply the baseline on the four conditions, with competition among them. Above the allotment, supply is priced at cost, including energy, and the market continues to exist where the market is useful. What disappears is the layer that charges the public rent for the continued existence of its own memory.

Two bodies of work bear on the design. Studies of long-lived common resources find that they tend to be governed through clearly defined boundaries, monitoring accountable to the users and graduated sanctions \cite{ostrom}, which matches the bounded allotment and the narrow, contestable sanction proposed here. And the digital preservation literature has long held that digital records are fragile because formats and media become obsolete, so that keeping a record means keeping the means of reading it \cite{rothenberg}, which is the ground of the portability and legibility conditions.

\section{Objections}

The first objection is that nothing is free, and that someone must pay. This is true and is not an objection. Someone pays for roads and libraries, and the question is who and by what mechanism. The proposal is that the public pays in common, in proportion to means, through the same fiscal arrangements that provide other infrastructure, and not individually and perpetually at the seller's rate. That the payer is a taxpayer who did not choose each item is true of every common good, and the answer is the political one, a public decision on a defined and bounded capability with the levy designed in proportion to means; the essay does not design it. It offers no budget, and a serious one would have to be built from operating data that the large providers do not publish; the figures in section~\ref{sec:rent} show only that the hardware price of a baseline allotment is small beside the rent charged for it. At the illustrative figure given there, a terabyte of replicated hardware for each of a million people would cost between \$72 million and \$123 million before power, buildings, staff and replacement, which gives an order of magnitude and no more.

The second objection is abuse. A free allotment invites misuse, such as storage of illicit material, mining, resale of capacity and automated sign-up for multiple allotments. The answer is that a quota is not a price. A cap per holder and metering above it limit consumption, mining and resale, and the existing apparatus of law handles illicit content. The harder problem is the identity needed to enforce a per-person cap, which raises its own privacy costs. A household or cooperative account without real-name registration is one way to reduce them, and the problem is acknowledged here as open.

The third objection is that public provision will be poor, that public information systems have often been slow and badly designed. It is a fair description of some history. The conditions in section~\ref{sec:freemeans} are meant to make it survivable: with open formats and portable identity, a bad operator can be abandoned without loss, and competition among paid operators can be sustained. The concentration the proposal targets is held by the sellers who now define the formats.

A harder objection concerns procurement. A public authority buying capacity at scale may find that only the largest firms can supply it, and the baseline could end by routing public money to the same sellers. The four conditions are the protection, since portability and legibility let the authority change operators and let smaller operators and cooperatives enter, but the risk is real and the essay does not claim to have removed it.

The fourth objection is that free baselines would reduce the incentive to build new services. The layers above the baseline remain open to commerce, and the commodity layer, bytes held and cycles supplied, is not where the innovations that the public values are found. Photographs, mail and search were built on top of storage and compute, and they do not become unbuildable when the underlying capacity is a utility. Nothing in the history of those services shows that they depended on the capacity beneath them being sold by the month. A bounded baseline leaves demand above it intact, and it widens the number of people able to build, since an allowance that does not lapse removes a cost that now stops small projects before they start.

The fifth objection is environmental. Free use without limit would waste energy, and the metering that now restrains use does at least send a price signal. The reply is that the baseline is bounded, that use above it is priced at cost including energy, and that a price signal set by a seller's margin is a poor proxy for the cost to the grid. A quota priced by energy is the better signal. Because the baseline is bounded per person, its energy draw can be planned for, whereas metered consumption has no ceiling other than what customers will pay.

The last concession belongs to the evidence. The sources used are drawn from price comparisons and trade analyses of varying rigor, one of them a vendor's press release, and none of them is a cost study of any provider. The argument that rent exceeds hardware by a wide margin is an argument about orders of magnitude and structure, and a full accounting of the cost of running storage at scale could narrow the gap. It would not close the argument, which rests as much on discretion, exit and dependence as on the arithmetic.

The argument also does not need the economics to come out in its favor. Suppose providers incur every cost they say they incur and earn only ordinary margins. A competitively priced and efficiently run service can still be an unsuitable foundation for a lifelong record, because continued access then depends on individual purchasing power, and the interruption of access can destroy the conditions of a person's intellectual and personal continuity. That is a reason for a society to fund a defined capability in common, in the way it funds other things whose absence would cut a person off from ordinary life. The price evidence in the earlier sections shows that the present arrangement is not forced by cost. The normative argument shows that it would be unsuitable even if it were.

\section{Conclusion}

The meter sits between a person and the means of keeping and transforming a record. It was installed when storage and compute were scarce, and it has survived their abundance by changing its justification. It now rests on the claim that someone must pay, which is true of every road in the country, and on the habit of treating a person's own archive as a commercial relationship with a landlord who may change the lease. The documented facts are modest and sufficient. Storage rented from the largest provider costs in a month about what the disk costs to buy. Retrieval is priced above storage, and a competitor shows that billing for it separately is a choice. Europe has scheduled the end of switching charges for January 2027, which treats the cost of leaving as a matter of competition and not of accounting. A free tier collected the world's photographs and later made new deposits count against a finite allowance. Quotas bill copies and revisions by their visible volume. And the price of hardware, on a plausible reading, moves with the orders of buyers who are not the public.

Capacity alone would not be enough. Unless the means of rebuilding it are published and taught, a free allowance would be only a better-furnished tenancy, and the dependence would survive the abolition of the fee. Against all this the proposal is simple: a baseline of storage and compute, free at the point of use, portable, not revocable and legible, funded as the infrastructure it has become. Everything above the baseline can be sold, and should be sold at what it costs. What should not be sold, by the month and at the provider's discretion, is the permission to remember.

\begin{thebibliography}{99}

\bibitem{filebase} Filebase. ``AWS S3 Pricing in 2026: What You'll Actually Pay.'' \url{filebase.com/blog/aws-s3-pricing-in-2026-what-youll-actually-pay/}. Reports S3 Standard at \$0.023 per GB-month for the first 50 TB in US East; the publisher sells competing S3-compatible storage.

\bibitem{hddhunt} HDDHunt, a commercial price-comparison site. ``HDD price per TB, August 2026.'' \url{hddhunt.com/blog/hdd-price-per-tb-august-2026/}. The page is titled for August 2026, but the price table it carries is a snapshot dated October 2026 of the cheapest new internal drives on Amazon US, roughly \$24 to \$41 per TB by capacity.

\bibitem{ecorpit} ecorpit.com. ``Cloud egress fees in 2026: what AWS, Azure, Google Cloud and Cloudflare R2 actually charge.'' \url{ecorpit.com/cloud-egress-fees-ai-inference-bill-comparison-2026/}. Egress rates verified June 2026; notes on the 2024 exit waivers.

\bibitem{cloudmagazin} cloudmagazin.com. ``EU Data Act: when cloud switching fees are abolished, what CIOs need to examine.'' 6 July 2026. \url{cloudmagazin.com/en/2026/07/06/eu-data-act-when-cloud-switching-fees-are-abolished-what-cios-need-to-examine/}

\bibitem{sciencealert} ScienceAlert. ``When Gmail was first announced, people thought it was an April Fools' joke.'' \url{sciencealert.com/when-gmail-was-first-announced-people-thought-it-was-an-april-fools-joke}. Gmail's one gigabyte against two or four megabytes at Hotmail and Yahoo; the article calls this about a hundred times, but the figures it gives imply a factor of several hundred.

\bibitem{gmailblog} Official Gmail Blog. ``Bringing it all together: 15 GB now shared between Drive, Gmail and Google+ Photos.'' May 2013. \url{gmail.googleblog.com/2013/05/bringing-it-all-together-15-gb-now.html}

\bibitem{dataact} Regulation (EU) 2023/2854 of the European Parliament and of the Council of 13 December 2023 on harmonised rules on fair access to and use of data (Data Act). Recitals 88, 97 and 99. \url{eur-lex.europa.eu/legal-content/EN/TXT/?uri=CELEX:32023R2854}

\bibitem{googlephotos} Google. Announcement of changes to Google Photos storage, 11 November 2020. \url{blog.google/products/photos/storage-changes/}

\bibitem{iphonecanada} iPhone in Canada. ``Google Photos is ending its free unlimited storage in June 2021.'' 11 November 2020. \url{iphoneincanada.ca/2020/11/11/google-photos-is-ending-its-free-unlimited-storage-in-june-2021/}

\bibitem{eracks} eRacks Open Source Systems. ``Hard drive prices up 40 to 90 percent since March: eRacks publishes its own six-month tracking.'' Press release, 14 September 2026. \url{pr.com/press-release/978414}. Vendor's own price tracking; its headline range does not match its table, which shows increases of 15 to 92 percent.

\bibitem{ostrom} Ostrom, E. (1990). \emph{Governing the Commons: The Evolution of Institutions for Collective Action}. Cambridge University Press.

\bibitem{rothenberg} Rothenberg, J. (1995). Ensuring the longevity of digital documents. \emph{Scientific American}, 272(1), 42--47.

\end{thebibliography}

\end{document}
